% Shared formatting for all notes in this repository.
% Adapted from the UC Berkeley EECS CS267 lecture-note template:
% https://www.cs.cmu.edu/~ggordon/10725-F12/template.tex
% Personal-note adaptation: boxed title/author/date; ordinary numbering;
% existing natbib citations and mathematical environments are retained.
% Load after \documentclass; compile from the repository root.
\usepackage[T1]{fontenc}
\usepackage[utf8]{inputenc}
\usepackage{amsmath,amsfonts,amssymb,mathrsfs,amsthm}
\usepackage{graphicx,wrapfig}
\usepackage[inline,shortlabels]{enumitem}
\usepackage{centernot}
\usepackage[numbers]{natbib}
\usepackage{adjustbox,booktabs,arydshln}
\usepackage{caption,subcaption}
\usepackage{xcolor,tikz,pifont}
\usetikzlibrary{decorations.pathreplacing,angles,quotes,shapes.geometric,
  arrows.meta,positioning,calc,fit,backgrounds,shadows}
\usepackage[ruled,vlined]{algorithm2e}
\usepackage{xr}
\usepackage{hyperref}
\hypersetup{colorlinks=true,linkcolor=blue,citecolor=blue,urlcolor=blue}

% Berkeley EECS lecture-note page layout and paragraph spacing.
\setlength{\oddsidemargin}{0.25in}
\setlength{\evensidemargin}{-0.25in}
\setlength{\topmargin}{-0.6in}
\setlength{\textwidth}{6.5in}
\setlength{\textheight}{8.5in}
\setlength{\headsep}{0.75in}
\setlength{\parindent}{0in}
\setlength{\parskip}{0.1in}

% A wrapping box accommodates long research-note titles.
\makeatletter
\renewcommand{\maketitle}{%
  \pagestyle{myheadings}%
  \thispagestyle{plain}%
  \begin{center}
    \fbox{%
      \begin{minipage}{\dimexpr\textwidth-2\fboxsep-2\fboxrule\relax}
        \vspace{2mm}
        \textbf{Research Notes}\hfill\@date\par
        \vspace{2mm}
        {\centering\Large\bfseries\@title\par}
        \vspace{2mm}
        {\itshape Author: \@author\par}
        \vspace{2mm}
      \end{minipage}%
    }
  \end{center}
  \markboth{Research Notes}{Research Notes}%
  {\footnotesize\itshape
    Style adapted from the UC Berkeley EECS lecture-note template.\par}
  \vspace{4mm}%
}
\makeatother

\newcommand{\bR}{{\boldsymbol R}}
\newcommand{\bX}{{\boldsymbol X}}
\newcommand{\bY}{{\boldsymbol Y}}
\newcommand{\by}{{\boldsymbol y}}
\newcommand{\bx}{{\boldsymbol x}} 
\newcommand{\bz}{{\boldsymbol z}}
\newcommand{\bZ}{{\boldsymbol Z}}
\newcommand{\bH}{{\boldsymbol H}}
\newcommand{\bb}{{\boldsymbol b}}
\newcommand{\be}{{\boldsymbol e}}
\newcommand{\bbf}{{\boldsymbol f}}
\newcommand{\bg}{{\boldsymbol g}}
\newcommand{\bm}{{\boldsymbol m}}
\newcommand{\bw}{{\boldsymbol w}}
\newcommand{\bW}{{\boldsymbol W}}
\newcommand{\bv}{{\boldsymbol v}}
\newcommand{\bt}{{\boldsymbol t}}
\newcommand{\bA}{{\boldsymbol A}}
\newcommand{\ba}{{\boldsymbol a}}
\newcommand{\bD}{{\boldsymbol D}}
\newcommand{\bE}{{\boldsymbol E}}
\newcommand{\bi}{{\boldsymbol i}}
\newcommand{\bc}{{\boldsymbol c}}
\newcommand{\bC}{{\boldsymbol C}}
\newcommand{\bk}{{\boldsymbol k}}
\newcommand{\bK}{{\boldsymbol K}}
\newcommand{\bP}{{\boldsymbol P}}
\newcommand{\mK}{\mathcal{K}}
\newcommand{\bV}{{\boldsymbol V}}
\newcommand{\bG}{{\boldsymbol G}}
\newcommand{\mG}{\mathcal{G}}
\newcommand{\mH}{\mathcal{H}}
\newcommand{\mE}{\mathbb{E}}
\newcommand{\mX}{\mathcal{X}}
\newcommand{\mM}{\mathcal{M}}
\newcommand{\bU}{{\boldsymbol U}}
\newcommand{\bu}{{\boldsymbol u}}
\newcommand{\bI}{{\boldsymbol I}}
\newcommand{\bJ}{{\boldsymbol J}}
\newcommand{\bs}{{\boldsymbol s}}

\newcommand{\bxi}{{\boldsymbol \xi}}
\newcommand{\balpha}{{\boldsymbol \alpha}}
\newcommand{\bbeta}{{\boldsymbol \beta}}
\newcommand{\bgamma}{{\boldsymbol \gamma}}
\newcommand{\btheta}{{\boldsymbol \theta}}
\newcommand{\bTheta}{{\boldsymbol \Theta}}
\newcommand{\bvartheta}{{\boldsymbol \vartheta}}
\newcommand{\bepsilon}{{\boldsymbol \epsilon}}
\newcommand{\bvarepsilon}{{\boldsymbol \varepsilon}}
\newcommand{\bomega}{{\boldsymbol \omega}}
\newcommand{\bmu}{{\boldsymbol \mu}}
\newcommand{\bSigma}{{\boldsymbol \Sigma}}

\newcommand{\corr}{\mbox{Corr}}
\newcommand{\var}{{\mbox{Var}}}
\newcommand{\Var}{{\mbox{Var}}}
\newcommand{\cov}{{\mbox{Cov}}}
\newcommand{\diag}{{\mbox{diag}}}
\newcommand{\bUmis}{{\boldsymbol U}_{{\rm mis}}}
\newcommand{\bYmis}{{\boldsymbol Y}_{{\rm mis}}}
\newcommand{\bYobs}{{\boldsymbol Y}_{{\rm obs}}}
\newcommand{\tbXmis}{{\tilde{\boldsymbol X}}_{{\rm  mis}}}
\newcommand{\tbxmis}{{\tilde{\boldsymbol x}}_{{\rm  mis}}}
\newcommand{\bXmis}{{\boldsymbol X}_{{\rm mis}}}
\newcommand{\bXobs}{{\boldsymbol X}_{{\rm obs}}}
\newcommand{\bZmis}{{\boldsymbol Z}_{{\rm mis}}}
\newcommand{\bZobs}{{\boldsymbol Z}_{{\rm obs}}}
\newcommand{\tXmis}{{\tilde{X}}_{{\rm  mis}}}
\newcommand{\txmis}{{\tilde{x}}_{{\rm  mis}}}
\newcommand{\Xmis}{X_{{\rm  mis}}}
\newcommand{\Xobs}{X_{{\rm  obs}}}
\newcommand{\bymis}{{\boldsymbol y}_{{\rm  mis}}}
\newcommand{\byobs}{{\boldsymbol y}_{{\rm  obs}}}
\newcommand{\ymis}{y_{{\rm  mis}}}
\newcommand{\yobs}{y_{{\rm  obs}}}

\newcommand{\JRSSB}{Journal of the Royal Statistical Society: Series B}
\newcommand{\indep}{\rotatebox[origin=c]{90}{$\models$}}
\newcommand{\red}{\textcolor{red}}
\newcommand{\blue}{\textcolor{blue}}
\newcommand{\notetodo}[1]{\textcolor{blue}{[TODO: #1]}}

% Preserve the existing independent, section-based theorem counters.
\newtheorem{theorem}{Theorem}[section]
\newtheorem{lemma}{Lemma}[section]
\newtheorem{corollary}{Corollary}[section]
\newtheorem{formula}{Formula}[section]
\newtheorem{proposition}{Proposition}[section]
\theoremstyle{definition}
\newtheorem{definition}{Definition}[section]
\newtheorem{example}{Example}[section]
\newtheorem{assumption}{Assumption}
\theoremstyle{remark}
\newtheorem{remark}{Remark}

% Journal abbreviations used by the existing bibliography.
\def\AJE{American Journal of Epidemiology}
\def\ANNALS{Annals of Statistics}
\def\BIOK{Biometrika}
\def\BMCS{Biometrics}
\def\STATMED{Statistics in Medicine}
\def\STATSCI{Statistical Science}
\def\JSPI{Journal of Statistical Planning \& Inference}
\def\JRSSB{Journal of the Royal Statistical Society, Series B}
\def\COMMS{Communications in Statistics, Series A}
\def\JQT{Journal of Quality Technology}
\def\TECH{Technometrics}
\def\JASA{Journal of the American Statistical Association}
\def\CSDA{Computational Statistics \& Data Analysis}
\def\JCGS{Journal of Computational and Graphical Statistics}
\def\JCB{Journal of Computational Biology}
\def\BIOINF{Bioinformatics}
\def\JAMA{Journal of the American Medical Association}
\def\JNUTR{Journal of Nutrition}
\def\LETTERS{Letters in Probability and Statistics}
\def\JABES{Journal of Agricultural, Biological and Environmental Statistics}
\def\SCAN{Scandinavian Journal of Statistics}
\def\PNAS{Proceedings of the National Academy of Sciences USA}
\def\JMLR{Journal of Machine Learning Research}
\def\JSCS{Journal of Statistical Computation and Simulation}
